% !TEX program = xelatex
\documentclass{misisarticle}
\usepackage{tikz}
\usepackage{needspace}
\usetikzlibrary{arrows.meta,positioning}
\grnti{}
\asjc{}
\oecd{}
\title{Исследование стратегий поэтапного обучения захвату и подъёму объектов роботом Unitree G1}
\misisauthor{Арепьев Егор Сергеевич}{НИТУ «МИСИС», группа БПМ-23-ПО-2}{Москва, Россия}{arepyev2006@gmail.com}
\misisauthor{Шепелев Д.\,А.}{Научный руководитель}{}{}
\abstract{Рассматривается обучение захвату и подъёму объекта с полки многопальцевой кистью робота Unitree G1 в симуляционной среде. Исследуются три стратегии: обучение общей политики для полной задачи, обучение отдельных навыков и обучение единой политики с последовательной сменой конечной цели. Описаны наблюдения, управление, функции награды и критерии устойчивого удержания. Предложен протокол сравнения по успешности и числу взаимодействий со средой. Приведены предварительные оценки в одной фиксированной сцене: после изменения наград и масштаба действий единая политика освоила подведение руки, однако успешность полной последовательности действий пока не подтверждена. Результаты показывают необходимость раздельно оценивать достижение контакта, подъём и устойчивость удержания.}
\keywords{обучение с подкреплением, роботизированный захват, PPO, поэтапное обучение, Unitree G1, Isaac Lab}
\begin{document}
\maketitle

\section{Введение}
Извлечение товара с заполненной полки требует согласованного управления рукой и кистью. Робот должен приблизиться к объекту, сформировать удерживающий контакт и поднять объект, не нарушив расположение соседних товаров. Успешное приближение ещё не означает успешного захвата, а кратковременный подъём не гарантирует устойчивого удержания. Для многопальцевой кисти дополнительные сложности возникают из-за изменения контактных сил и вращения объекта внутри хвата.

Обучение с подкреплением позволяет оптимизировать последовательность действий по сигналу награды. Однако при вознаграждении только за полное выполнение задачи случайная политика может редко достигать полезных состояний. Промежуточные награды и упрощённые начальные условия облегчают исследование среды, но могут закреплять нежелательное поведение, например касание объекта без последующего подъёма.

Цель работы --- исследовать стратегии поэтапного обучения захвату и подъёму товара роботом Unitree G1 с кистью Inspire в сцене RoboBenchMart. Рассматриваются одна политика с неизменной конечной задачей, отдельная политика для каждого навыка и одна политика с последовательной сменой конечной цели. Основной вопрос состоит в том, как эти варианты влияют на освоение промежуточных действий и устойчивого удержания при ограниченном бюджете взаимодействий со средой.

Работа использует существующие модели робота, сцену и базовую инфраструктуру обучения. Исследуемый вклад состоит в реализации вариантов декомпозиции задачи, разработке наград и критериев этапов, а также диагностике результатов. Здесь представлены реализованные методы и предварительные наблюдения; завершённое сравнение полной задачи по единому протоколу остаётся следующим этапом исследования.

\section{Анализ работ}
\subsection{Обучение политик управления}
Алгоритм Proximal Policy Optimization (PPO) чередует сбор траекторий текущей политикой и несколько проходов оптимизации по собранным данным~\cite{ppo}. Ограниченная целевая функция уменьшает стимул к чрезмерному изменению отношения вероятностей действий. В настоящей работе PPO используется как базовый алгоритм; предмет исследования --- организация задач и сигналов награды, а не разработка нового оптимизатора.

В работе OpenAI и соавторов показано обучение переориентации объекта многопальцевой кистью в симуляции с последующим переносом на физического робота~\cite{dexterous}. Существенную роль играет рандомизация физических параметров. Рассматриваемая здесь задача отличается наличием полки и соседних объектов: требуется сначала сформировать хват, а затем вывести объект из опоры. В текущих экспериментах параметры фиксированы, а перенос на реального робота не проверяется.

\subsection{Поэтапное обучение}
Narvekar и соавторы систематизируют подходы, в которых задачи или обучающие данные предъявляются в определённой последовательности~\cite{curriculum}. Для захвата необходимо различать изменение начальных состояний, изменение награды и изменение конечного критерия эпизода. Эти механизмы могут применяться совместно, но отвечают на разные исследовательские вопросы.

Florensa и соавторы предлагают начинать обучение вблизи целевого состояния, постепенно расширяя множество стартов~\cite{reverse}. Это мотивирует использование облегчённых стартов при освоении отдельных навыков. В наших экспериментах близкие к хвату старты задаются вручную и постепенно исключаются; алгоритм автоматической генерации стартов из указанной работы не воспроизводится. В варианте единой политики со сменой целей эпизоды начинаются из исходной позы руки.

\subsection{Симуляционная среда}
RoboBenchMart предназначен для исследования роботизированных операций в торговых средах с плотным размещением товаров~\cite{robobenchmart}. Здесь используется фиксированная сцена на основе этого проекта, адаптированная в рабочей инфраструктуре для Isaac Lab и G1. Это отдельный эксперимент с управлением рукой, а не повторение полного бенчмарка мобильного манипулирования.

Isaac Lab предоставляет инфраструктуру GPU-параллельной симуляции и обучения роботизированных политик~\cite{isaaclab}. Параллельные экземпляры среды позволяют собирать большой объём переходов, однако сами по себе не устраняют неудачный дизайн награды. Поэтому стратегии следует сравнивать не только по числу обновлений PPO, но и по объёму взаимодействий со средой.

\section{Постановка задачи}
\subsection{Сцена и управление}
Используется Unitree G1 с правой кистью Inspire. База робота зафиксирована; навигация и поддержание равновесия всего тела в задачу не входят. Целевой объект --- банка Nutella среди соседних товаров на полке. Масса банки составляет $0{,}61$~кг; коэффициенты трения банки, робота и полки равны соответственно $0{,}5$, $1{,}1$ и $0{,}4$. Это принятые параметры модели, а не результат экспериментальной калибровки реального контакта. Используются сцепка фаланг и ограничение момента $1$~Н$\cdot$м на управляемый сустав кисти.

Физика рассчитывается с частотой $200$~Гц, действия задаются с частотой $50$~Гц. Полная политика выдаёт 19 компонент действия: 7 для руки и 12 для суставов кисти. Действие изменяет целевое положение суставов PD-регулятора:
\begin{equation}
 q^{\mathrm{cmd}}_{t+1}=\operatorname{clip}
 \bigl(q^{\mathrm{cmd}}_t+S a_t,\,q_{\min},q_{\max}\bigr),
\end{equation}
где $a_t\in[-1,1]^{19}$, а $S$ задаёт масштабы приращений. Компоненты кисти дополнительно согласуются с моделью сцепки фаланг.

Наблюдения включают положения и скорости суставов, состояние объекта, геометрию и силы контактов, сведения о соседях и предыдущее действие. Это признаки из симуляции; распознавание объектов по изображению здесь не обучается. В единой политике со сменой целей дополнительно используются ошибки относительно поз подхода и хвата, разность команды и положения суставов и код конечной цели. Размер входа актора равен 182, критика --- 194 с учётом 12 дополнительных признаков. Для временного контекста используются двухслойные LSTM с размером скрытого состояния 256.

\subsection{Этапы и критерии успеха}
Полная задача разделена на подведение открытой кисти (Approach), формирование хвата (Grasp) и подъём с удержанием (Lift). В табл.~\ref{tab:gates} приведены критерии единой политики со сменой целей. Условия должны выполняться непрерывно в течение заданного времени; нарушение сбрасывает таймер. Падение объекта и нарушение допустимого состояния соседей считаются неуспехом и имеют приоритет над достижением цели. Таймаут не является успехом.

\begin{table}[tb]
\caption{Критерии этапов единой политики}
\label{tab:gates}
\centering\small
\begin{tabular}{@{}p{0.18\columnwidth}p{0.76\columnwidth}@{}}
\toprule
Этап & Условия \\
\midrule
Approach & Максимальная ошибка суставов руки $\leq0{,}08$ рад; ошибка открытой кисти $\leq0{,}15$ рад; скорости суставов $\leq0{,}25$ рад/с, ладони $\leq0{,}06$ м/с и $\leq0{,}5$ рад/с; сдвиг банки $\leq1{,}5$ см. Удержание $0{,}3$~с. \\
\addlinespace
Grasp & Пять допустимых контактов силой $\geq1$ Н на палец и геометрическое противопоставление; скорость банки $\leq0{,}08$ м/с и $\leq1$ рад/с; наклон $\leq0{,}35$ рад, сдвиг $\leq2{,}5$ см. Удержание $0{,}3$~с. \\
\addlinespace
Lift & После Approach и Grasp: подъём нижней точки на $\geq5$ см, сохранение пяти контактов; опорные силы от полки и соседей $<0{,}5$ Н; те же ограничения скоростей и наклона. Удержание $1$~с. \\
\bottomrule
\end{tabular}
\end{table}

Контакты ниже допустимой области объекта не засчитываются; проверяются расположение контактных точек по разные стороны объекта и достаточный размах. Контакт руки с полкой исключает готовность этапа. Отдельные ранее обученные Grasp и Lift используют более мягкое требование: большой палец и ещё минимум два. У общей исходной политики также отличаются высота и длительность удержания. Поэтому их текущие значения успешности не являются результатами единого сравнительного теста.

\section{Описание методов}
\subsection{Три стратегии обучения}
\textit{Общая политика с одной конечной задачей.} Одна модель обучается всей последовательности захвата и подъёма. В последовательных запусках могут изменяться награды и помощь в начальных состояниях, но конечная задача остаётся прежней. Промежуточный контакт сам по себе не завершает эпизод.

\textit{Отдельные политики.} Approach, Grasp и Lift обучаются независимо с нуля в специализированных средах. Для Approach управляются только семь суставов руки, а кисть принудительно поддерживается открытой. Grasp начинается из предзахвата; часть обучающих стартов приближена к контактной позе, и доля помощи уменьшается. Lift начинается из банка физических состояний готового хвата. Использованный банк содержит 19 состояний от ранее обученной общей политики. Поэтому эта проверка Lift ещё не проверяет передачу состояния от нового Grasp. Полная цепочка переключения трёх моделей пока не валидирована.

\Needspace{4\baselineskip}
\textit{Единая политика со сменой конечной цели.} Используется одна модель с неизменными интерфейсом и архитектурой. Сначала эпизод завершается успешным Approach, затем конечной целью становится Grasp, после этого --- Lift. На поздних этапах предыдущие достижения регистрируются внутри эпизода без сброса физического состояния и памяти LSTM. Между запусками сохраняются веса, оптимизатор и нормализация; состояния симулятора и память эпизодов инициализируются заново. Рука стартует из исходной позы с открытой кистью, но все 19 компонент действия определяются политикой.

\begin{figure}[tb]
\centering
% Original vector diagram; no simulator assets or screenshots.
\begin{tikzpicture}[
 font=\small,
 box/.style={draw=black!65,rounded corners=2pt,align=center,
 text width=2.1cm,minimum height=0.75cm,inner sep=3pt},
 flow/.style={-{Latex[length=2mm]},thick},
 title/.style={anchor=west,font=\small\bfseries}]
\node[title] at (0,0) {Одна конечная задача};
\node[box,text width=7.3cm] at (3.7,-.65)
 {Одна политика: подвести руку, схватить, поднять};
\node[title] at (0,-1.45) {Отдельные модели};
\node[box] (a) at (1.1,-2.15) {$\pi_A$: Approach};
\node[box] (g) at (3.7,-2.15) {$\pi_G$: Grasp};
\node[box] (l) at (6.3,-2.15) {$\pi_L$: Lift};
\draw[flow,dashed] (a)--(g);
\draw[flow,dashed] (g)--(l);
\node[title] at (0,-2.95) {Одна модель, последовательность целей};
\node[box] (ca) at (1.1,-3.65) {$\pi_\theta$ при цели A};
\node[box] (cg) at (3.7,-3.65) {$\pi_\theta$ при цели G};
\node[box] (cl) at (6.3,-3.65) {$\pi_\theta$ при цели L};
\draw[flow] (ca)--(cg);
\draw[flow] (cg)--(cl);
\end{tikzpicture}

\caption{Организация трёх стратегий. Нижние стрелки обозначают перенос параметров одной модели между этапами обучения. Цепочка отдельных навыков показана как целевая схема интеграции, её сквозная оценка ещё не выполнена.}
\label{fig:strategies}
\end{figure}

Для новой единой политики используются 2000 сред и 96 шагов сбора на обновление PPO. Каждые 50 обновлений выполняются две детерминированные оценки по 100 эпизодов с разными seed начальных состояний. Переход разрешается, если успешность в каждой оценке не ниже $90\%$. Это порог продвижения, а не требование к итоговому качеству. Бюджет ограничен 500 обновлениями на этап. Две оценки одной модели не заменяют повторное обучение с независимыми seed.

\subsection{Оптимизация политики}
Используется ограниченная целевая функция PPO~\cite{ppo}:
\begin{equation}
\begin{split}
 L^{\mathrm{clip}}(\theta)=\mathbb{E}_t\bigl[\min\{&\rho_t(\theta)\hat A_t,\\
 &\operatorname{clip}(\rho_t(\theta),1-\epsilon,1+\epsilon)\hat A_t\}\bigr],
\end{split}
\end{equation}
где $\rho_t=\pi_\theta(a_t\mid o_t,h_t)/\pi_{\mathrm{old}}(a_t\mid o_t,h_t)$, $h_t$ --- рекуррентное состояние, $\hat A_t$ --- оценка преимущества. Для новой единой политики приняты $\epsilon=0{,}2$, $\gamma=0{,}99$, $\lambda=0{,}95$, начальная скорость обучения $3\cdot10^{-4}$, 5 эпох оптимизации и 8 мини-пакетов. Скорость обучения адаптируется по расхождению политик. Эти параметры сохранены при изменении награды.

\subsection{Награды и помощь в обучении}
Первоначальный вариант единой политики поощрял приближение разностью потенциалов. Для нетерминального перехода при $\Phi=-2e$, где $e$ --- средняя ошибка суставов, уменьшение $e$ с $0{,}30$ до $0{,}29$ рад давало $0{,}0258$ награды. Рабочий отдельный Approach за такое уменьшение получал компонент прогресса $+1$ и штраф за оставшееся расстояние. Кроме того, его масштаб приращений руки был в пять раз больше. Это мотивировало проверку более выраженного сигнала достижения позы, но не доказывает причинную роль отдельного коэффициента.

В новой версии основные компоненты награды подхода имеют вид
\begin{equation}
\begin{split}
 r^{\mathrm{reach}}_t={}&2\operatorname{clip}\!\left(\frac{e_{t-1}-e_t}{0{,}02},-1,1\right)\\
 &-2e_t-0{,}5e_t^{\max}.
\end{split}
\end{equation}
Добавлены штрафы за закрывание кисти и избыточную команду сгибания, небольшие бонусы за спокойную позу вблизи цели и готовность. Компоненты подхода сохраняются до достижения Approach и на поздних этапах. Приращения плеча и остальных суставов руки увеличены с $0{,}006/0{,}008$ до $0{,}015/0{,}02$ рад на единицу действия; масштаб кисти оставлен $0{,}04$ рад.

Для Grasp введены награды за число контактирующих пальцев, большой палец, противопоставление и полный пятипальцевый хват. Положительные компоненты частичного контакта выплачиваются не более 80 контактных шагов за эпизод. Бонусы за новый максимум числа пальцев и конфигурации хвата выдаются однократно. Это ограничивает получение награды за повторение касаний без завершения задачи.

Для Lift награда за высоту отделена от награды за спокойное удержание. При допустимом хвате, отсутствии опоры и безопасном состоянии сцены соответствующие компоненты имеют вид
\begin{equation}
 \begin{split}
 r^{\mathrm{height+calm}}_t={}&0{,}35H_t+0{,}15\\
 &+\frac{0{,}4H_t}{1+(v_t/0{,}08)^2+(\omega_t/1)^2}.
 \end{split}
\end{equation}
где $H_t=\operatorname{clip}(h_t/0{,}05,0,1)$, $h_t$ --- подъём нижней точки объекта в метрах, $v_t$ и $\omega_t$ --- модули скоростей в м/с и рад/с. Компоненты активируются при подъёме не менее 5 мм. Награда за новый максимум высоты ограничена бюджетом до достижения 5 см; повторный подъём на уже достигнутую высоту не оплачивается. Формула описывает часть полной награды: дополнительно учитываются хват, готовность, резкость действий и безопасность.

Бонусы за конечный успех Approach, Grasp и Lift равны 100, 150 и 200. Положительные награды не отменяют ограничений безопасности. Строгие критерии оценки сохранены; физические параметры, внешние силы и демпфирование для облегчения задачи не менялись. Новая политика обучается с нуля. Совместное изменение наград и масштаба управления является проверкой пакета изменений, а не изолированной абляцией.

\Needspace{11\baselineskip}
\section{Предварительные результаты и протокол оценки}
\subsection{Подтверждённые наблюдения}
В табл.~\ref{tab:results} приведён срез завершённых оценок, проверенных 6 ноября 2026 года. Неполные журналы не используются как результаты. Успешность вычисляется как
\begin{equation}
 \mathrm{SR}=\frac{1}{N}\sum_{i=1}^{N}\mathbf{1}\{\text{эпизод }i\text{ завершён успехом}\}.
\end{equation}

\begin{table}[tb]
\caption{Предварительные оценки: разные задачи и условия старта}
\label{tab:results}
\centering\small
\begin{tabular}{@{}p{0.56\columnwidth}r@{}}
\toprule
Модель и оцениваемый этап & Успехи / эпизоды \\
\midrule
Общая политика, строгий подъём & $0/64$ \\
Отдельный Approach & $16/16$ \\
Отдельный Grasp, checkpoint 175 & $62/64$ \\
Отдельный Lift, checkpoint 999 & $0/57$ \\
Единая, прежний Approach & $0/100$; $0/100$ \\
Единая, новый Approach & $100/100$; $100/100$ \\
Единая, новый Grasp & $0/100$; $0/100$ \\
\bottomrule
\end{tabular}
\end{table}

У общей политики факт удержания наблюдался в 64 из 64 эпизодов, но строгий критерий завершения не выполнялся. Отдельный Grasp оценён на промежуточном checkpoint 175; результат $62/64$ не приписывается финальному checkpoint 999. Для отдельного Lift 55 из 57 эпизодов достигали высоты 5 см с допустимым хватом и отсутствием опоры, но полного строгого успеха не было. Эти 57 эпизодов --- три повторения 19 банковых состояний, а не 57 независимых сцен.

Первоначальная единая политика не освоила Approach за 500 обновлений, то есть 96 млн переходов. Дополнительная оценка из 16 эпизодов показала среднюю максимальную ошибку суставов в конце эпизода $0{,}83$ рад при допуске $0{,}08$ рад. Открытие кисти и ограничения скоростей в этих конечных состояниях выполнялись; препятствием было недостижение целевой позы.

После изменения наград и масштаба действий новая единая политика достигла $100/100$ в каждой из двух оценок Approach после 50 обновлений, или 9,6 млн переходов. Средняя максимальная ошибка суставов в момент успеха составляла около $0{,}022$ рад. Эти величины описывают разные моменты завершения эпизодов и не являются оценкой ошибки на одном фиксированном шаге. После перехода к Grasp последние завершённые оценки на 300 обновлениях этого этапа дали $0/100$ в обоих прогонах; достижение Approach сохранялось во всех эпизодах. Подтверждено освоение первого этапа, но не всей последовательности.

\subsection{Ограничения и дальнейшее сравнение}
Число обновлений недостаточно для оценки стоимости обучения. При постоянном числе сред $n$ и длине сбора $T$ бюджет равен
\begin{equation}
 B=nTK,
\end{equation}
где $K$ --- число обновлений. Отдельный Approach обучался 250 обновлений при $n=128$, $T=32$, то есть использовал 1,024 млн переходов. Новый единый Approach использовал 9,6 млн. Меньшее число обновлений во втором случае не означает меньших затрат данных; различаются также размер действия и управление пальцами.

Для итогового сравнения планируется общий набор начальных состояний, единый критерий полного подъёма, одинаковая физика и несколько независимых обучений. Помимо SR следует измерять доли достижения хвата и высоты, максимальную длительность непрерывного удержания, частоту падений и нарушений состояния соседей. Бюджет отдельных навыков необходимо суммировать, а стоимость подготовки банка начальных хватов учитывать отдельно.

Оценки относятся к одной сцене и одному объекту с ограниченным шумом начальных состояний. Они не подтверждают обобщение на новые товары, успешность порядка $99\%$ в широком распределении условий или перенос на физического робота. IK и cuRobo рассматриваются как последующее расширение, но не включены в представленные обученные методы. Следующая задача --- освоение перехода от Approach к Grasp и устойчивого Lift, после чего станет возможна сквозная оценка трёх стратегий.

\begin{thebibliography}{9}
\bibitem{ppo} Schulman J., Wolski F., Dhariwal P., Radford A., Klimov O. Proximal Policy Optimization Algorithms. arXiv:1707.06347, 2017. \url{https://arxiv.org/abs/1707.06347}.
\bibitem{dexterous} OpenAI et al. Learning Dexterous In-Hand Manipulation. arXiv:1808.00177, 2018. \url{https://arxiv.org/abs/1808.00177}.
\bibitem{curriculum} Narvekar S., Peng B., Leonetti M., Sinapov J., Taylor M. E., Stone P. Curriculum Learning for Reinforcement Learning Domains: A Framework and Survey. \textit{Journal of Machine Learning Research}, vol.~21, no.~181, pp.~1--50, 2020. \url{https://arxiv.org/abs/2003.04960}.
\bibitem{reverse} Florensa C., Held D., Wulfmeier M., Zhang M., Abbeel P. Reverse Curriculum Generation for Reinforcement Learning. \textit{Conference on Robot Learning}, 2017. arXiv:1707.05300. \url{https://arxiv.org/abs/1707.05300}.
\bibitem{robobenchmart} Soshin K., Krapukhin A., Spiridonov A., Bukhtuev G., Kuznetsov A., Shakhuro V., Shepelev D. RoboBenchMart: Benchmarking Robots in Retail Environment. arXiv:2511.10276, 2025. \url{https://arxiv.org/abs/2511.10276}.
\bibitem{isaaclab} Mittal M., Roth P., Tigue J. et al. Isaac Lab: A GPU-Accelerated Simulation Framework for Multi-Modal Robot Learning. arXiv:2511.04831, 2025. \url{https://arxiv.org/abs/2511.04831}.
\end{thebibliography}
\end{document}
